\section{Future Work}\label{sec:future}

\subsection{Immediate: Identify the Conserved Quantity Q}

Our experiments show that the system maintains a conserved quantity with extraordinary precision (EDI = 0.611991643906, 12 decimals, 5 seeds), but we have not identified what Q is mechanistically. Candidates include:

\begin{itemize}
\item \textbf{Mutual information} between layers or between attention heads
\item \textbf{Effective rank} of attention weight matrices
\item \textbf{Attention budget allocation} across heads (sum of attention entropies)
\item \textbf{Gradient norm ratios} between fast and slow components
\item \textbf{Representational capacity} measured via intrinsic dimensionality
\end{itemize}

\textbf{Proposed experiment:} Instrument training runs to log all candidate metrics (MI, Effective Rank, Attention Budget) at every step. We specifically propose to measure the cross-correlation between these candidates and the 12-decimal EDI plateau observed in Experiment 1B.

\paragraph{Preliminary Confirmation: Q is Effective Rank.}
We conducted a pilot analysis of 5 random seeds (0, 99, 101, 202, 303) from Experiment 1B. Remarkably, \textbf{all 5 seeds converged to an Effective Rank (ER) of exactly 4.267}. This value was maintained stably for thousands of steps (Table~\ref{tab:q-verification}). This precise cross-seed convergence is strong evidence that $Q$ is mechanistically the Effective Rank of the attention weight matrices.

\begin{table}[h]
\centering
\caption{Validation of Q Candidate. Effective Rank (ER) converges to 4.267 across all 5 tested seeds.}
\label{tab:q-verification}
\begin{tabular}{lcl}
\toprule
Seed & Step Locked & Final ER \\
\midrule
0 & 5,500 & \textbf{4.267} \\
99 & 500 & \textbf{4.267} \\
101 & 1,000 & \textbf{4.267} \\
202 & 1,000 & \textbf{4.267} \\
303 & 5,000 & \textbf{4.267} \\
\bottomrule
\end{tabular}
\end{table}

\subsection{Near-Term Experiments}

\paragraph{Timescale Ablation: Why 0.44--0.46?}
The forced attractor sits at EDI $\approx$ 0.44--0.46 under dual-timescale training with slow\_lr = 0.1. Does this value depend on the timescale separation?

\textbf{Design:} Run intentional\_edi\_target with target = 0.65 (to test ceiling) across slow\_lr $\in$ \{0.01, 0.05, 0.1, 0.5, 1.0\}.

\textbf{Hypothesis:}
\begin{enumerate}
\item \textbf{Soft constraint:} Forced attractor moves with slow\_lr (e.g., slower regulation allows higher EDI)
\item \textbf{Hard constraint:} Forced attractor stays at 0.44--0.46 regardless of timescale (geometric limit)
\end{enumerate}

If (1), the constraint is from dual-timescale dynamics. If (2), the constraint is from information geometry itself.

\paragraph{Lambda Sweep: Is It Convergence Loss or Timescale?}
Is the forced attractor caused by dual-timescale separation or by the convergence loss ($\lambda_{\text{conv}}$)?

\textbf{Design:} Fix slow\_lr = 0.1, vary $\lambda_{\text{conv}} \in$ \{0.0, 0.1, 0.3, 0.5, 1.0\}, target EDI = 0.65.

\textbf{Expected outcome:} Reveals whether the ceiling is from regulatory pressure (convergence loss penalizing head disagreement) or from timescale separation itself.

\paragraph{Kill Test on Real Suppressors (GPT-2, Mistral).}
Experiment 3 validated Le Chatelier dynamics on a grokking task. Does it generalize to real suppressor heads in large language models?

\textbf{Design:}
\begin{enumerate}
\item Fine-tune GPT-2 Medium on a factual recall task
\item Perturb heads 0:2, 0:4, 0:7 by scaling weights $\omega = 0.5$
\item Three conditions:
  \begin{itemize}
  \item Free: All other heads trainable (compensation allowed)
  \item Frozen: Layer-0 heads 0, 1, 5, 6 frozen (compensation blocked)
  \item Baseline: No perturbation
  \end{itemize}
\item Measure convergence speed and final calibration
\end{enumerate}

\textbf{Prediction:} Frozen condition will show slower convergence or degraded calibration, replicating kill test findings.

\subsection{Long-Term Directions}

\paragraph{Homeostatic Alignment.}
Can we design training regimes where safe behavior is the natural equilibrium?

\textbf{Steps:}
\begin{enumerate}
\item Measure VDI (or Q) for safety-relevant circuits in a baseline model
\item Identify where the natural equilibrium sits (suppressor-dominated or amplifier-dominated)
\item Design dual-timescale + homeostatic loss targeting a safe equilibrium
\item Validate via kill tests: does blocking compensation collapse alignment?
\end{enumerate}

\textbf{Success criterion:} Aligned behavior persists even when you perturb the training objective, because it is actively maintained as a homeostatic set point.

\paragraph{Predicting Equilibria from Task Structure.}
Can we predict where the natural equilibrium will settle given task geometry, before training?

\textbf{Approach:}
\begin{enumerate}
\item Characterize tasks by information-theoretic properties (entropy, MI, compression ratio)
\item Train models on multiple tasks, measure natural VDI equilibrium for each
\item Build a predictive model: task properties → equilibrium location
\item Test on held-out tasks
\end{enumerate}

\textbf{Impact:} If successful, this would allow principled design of training curricula to shape equilibrium structure predictably.

\paragraph{Cross-Architecture Generalization.}
Do forced attractors exist in non-transformer architectures?

\textbf{Candidates:} CNNs, RNNs, state-space models (Mamba, S4)

\textbf{Design:} Apply dual-timescale training + VDI target sweep to each architecture. Measure whether forced attractors emerge and where they sit.

\textbf{Hypothesis:} Forced attractors are a general property of dual-timescale training under information bottleneck constraints, not specific to transformers.

\paragraph{Information-Geometric Analysis of Forced Attractors.}
Why does the forced attractor sit at VDI $\approx$ 0.44--0.46?

\textbf{Proposed measurements:}
\begin{itemize}
\item Effective rank of attention matrices at different VDI targets
\item Mutual information between layers throughout training
\item Attention entropy distributions (is 0.44 a special value?)
\item Gradient flow analysis (where do gradients saturate?)
\end{itemize}

\textbf{Goal:} Provide a mechanistic explanation for the forced attractor location, connecting it to measurable geometric properties.

\subsection{Theoretical Directions}

\paragraph{Formalize Homeostatic Manifolds.}
Develop a mathematical framework for equilibrium manifolds in neural networks:
\begin{itemize}
\item Define the manifold $\mathcal{M}$ as the set of states where $|\frac{dQ}{dt}| < \epsilon$
\item Characterize boundaries as regions where $\frac{\partial Q}{\partial \theta}$ diverges (saturation)
\item Prove existence of forced attractors under dual-timescale gradient flow
\item Derive conditions for phase transitions between equilibria
\end{itemize}

\paragraph{Connect to Thermodynamics.}
Le Chatelier's principle is a thermodynamic law. Can we formalize the analogy?
\begin{itemize}
\item Define an analog of "free energy" for neural networks (loss + entropy term?)
\item Show that compensation minimizes this free energy
\item Prove that forced attractors are local minima of free energy under constraints
\item Connect to maximum entropy principles in statistical physics
\end{itemize}

\paragraph{Generalize to Continual Learning.}
Homeostatic equilibria may explain catastrophic forgetting:
\begin{itemize}
\item Old tasks settle at one equilibrium (VDI = $Q_1$)
\item New tasks push the system toward a different equilibrium (VDI = $Q_2$)
\item If $Q_1$ and $Q_2$ are incompatible (forced attractors don't overlap), forgetting occurs
\item Solution: Design training to create overlapping attractor basins
\end{itemize}
